% Gesamttest «Texte, Geschwindigkeit und Begegnung» — Quelle des PDFs (python3 scripts/build-lp-pdf.py bewegung).
% Musterlösung und Raster: bewertungspaket.tex. Alle Werte mit python3 nachgerechnet (scripts/lp/bewegung/zahlen.py).
% Vorgabe ~/TODO-BM-Vorkurs-LP03-15.md §5: Gleichung aus Text mit Variable und Antwortsatz (G1, G2), s = v · t und Einheiten
% und Plausibilität (G3, G4), Abschnitte tabellarisch und am Graphen (G5, G6), Begegnung mit Modellannahme und Einholen
% (G7, G8). 25 Punkte, rund 30 Minuten.
\documentclass[11pt]{article}
\usepackage{../lp-druck}
\renewcommand{\lpname}{Texte, Geschwindigkeit und Begegnung}
\renewcommand{\lpteilgebiet}{}
\renewcommand{\lpfuss}{Gesamttest · 25 Punkte}
\begin{document}
\lpkopf{Gesamttest}{%
  \vspace{4pt}\noindent\begin{tabularx}{\linewidth}{@{}X X@{}}
  Code (kein Name): \hrulefill & Punkte: \hrulefill\ / 25
  \end{tabularx}}

\begin{lpinfo}{Anleitung}
\textbf{Zeit:} rund 30 Minuten. \textbf{Hilfsmittel:} keine, kein Taschenrechner.
Schreib zu jeder Aufgabe den \textbf{Rechenweg} mit Einheiten und, wo eine Frage gestellt ist, einen \textbf{Antwortsatz} auf; bewertet wird der Weg.
Danach: selbst mit dem \emph{Bewertungspaket} bewerten oder die Lösung samt diesem Blatt scannen bzw.\ fotografieren
(ohne Namen und Standort) und mit dem Bewertungspaket einer KI geben.
\end{lpinfo}

\lpteil{Teil A · Vom Text zur Gleichung}{Kapitel 1 · 7 Punkte}

\begin{lpaufgabe}{G1}{Streaming}{4 P}
Anbieter A verlangt $9$ Fr. im Monat und $3$ Fr. pro Film. Anbieter B verlangt $30$ Fr. im Monat, Filme inbegriffen.
Ab wie vielen Filmen im Monat ist B günstiger? Lege eine Variable fest, stelle eine Gleichung auf, löse sie, mach die Probe
und schreib einen Antwortsatz.
\lpplatz{7}
\end{lpaufgabe}

\begin{lpaufgabe}{G2}{Zahlenrätsel}{3 P}
Drei aufeinanderfolgende \textbf{gerade} Zahlen haben zusammen die Summe $84$. Welche Zahlen sind es?
Schreib die Gleichung auf und mach die Probe.
\lpplatz{4}
\end{lpaufgabe}

\lpteil{Teil B · Weg, Geschwindigkeit, Zeit}{Kapitel 2 · 6 Punkte}

\begin{lpaufgabe}{G3}{Zugfahrt}{3 P}
Ein Zug fährt $135$ km in $1$ h $30$ min.
\begin{enumerate}[label=(\alph*),leftmargin=*,itemsep=3pt]
\item Wie gross ist seine Durchschnittsgeschwindigkeit?
\item Wie lange braucht er mit dieser Geschwindigkeit für $60$ km? Gib die Zeit in Minuten an.
\end{enumerate}
\lpplatz{4}
\end{lpaufgabe}

\begin{lpaufgabe}{G4}{Einheiten}{3 P}
Rechne um und schreib den Rechenweg auf.
\begin{enumerate}[label=(\alph*),leftmargin=*,itemsep=3pt]
\item $54$ km/h in m/s
\item $36$ min in Stunden (als Dezimalzahl)
\item Ein Velo fährt $5$ m/s. Wie viele km/h sind das? Ist das plausibel?
\end{enumerate}
\lpplatz{4}
\end{lpaufgabe}

\lpteil{Teil C · Fahrt in Abschnitten}{Kapitel 3 · 6 Punkte}

\begin{lpaufgabe}{G5}{Bergwanderung}{4 P}
Eine Wanderung: $9$ km mit $4.5$ km/h, dann $30$ min Pause, dann $3$ km steil bergauf mit $2$ km/h.
\begin{enumerate}[label=(\alph*),leftmargin=*,itemsep=3pt]
\item Erfasse die drei Abschnitte in einer Tabelle mit Geschwindigkeit, Zeit und Weg.
\item Wie lange dauert die ganze Wanderung?
\item Wie gross ist die Durchschnittsgeschwindigkeit der ganzen Wanderung?
\item Nina sagt: «Im Schnitt $3.25$ km/h, das Mittel von $4.5$ und $2$.» Erkläre, warum das nicht stimmt.
\end{enumerate}
\lpplatz{9}
\end{lpaufgabe}

\begin{lpaufgabe}{G6}{Graph lesen}{2 P}
Der Graph zeigt eine Velofahrt ($t$ in h, $s$ in km).
\begin{center}
\begin{tikzpicture}[x=16mm,y=1.6mm,font=\scriptsize]
  \draw[lplinie,line width=0.3pt,xstep=0.5,ystep=6] (0,0) grid (3,30);
  \draw[->,line width=0.6pt] (0,0) -- (3.3,0) node[right] {$t$ [h]};
  \draw[->,line width=0.6pt] (0,0) -- (0,33) node[above] {$s$ [km]};
  \foreach \i in {0.5,1,1.5,2,2.5,3} \node[below,inner sep=1.5pt] at (\i,0) {$\i$};
  \foreach \j in {6,12,18,24,30} \node[left,inner sep=3pt] at (0,\j) {$\j$};
  \draw[lpblau,line width=1.2pt] (0,0) -- (1,18) -- (1.5,18) -- (3,30);
\end{tikzpicture}
\end{center}
\begin{enumerate}[label=(\alph*),leftmargin=*,itemsep=3pt]
\item In welchem Abschnitt fährt die Person am schnellsten? Woran siehst du das im Graphen?
\item Wie gross ist die Durchschnittsgeschwindigkeit der ganzen Fahrt?
\end{enumerate}
\lpplatz{4}
\end{lpaufgabe}

\lpteil{Teil D · Begegnung und Einholen}{Kapitel 4 · 6 Punkte}

\begin{lpaufgabe}{G7}{Begegnung}{3 P}
Zwei Dörfer liegen $50$ km auseinander. Gleichzeitig fährt Tom von Dorf A mit $18$ km/h und Sara von Dorf B mit $12$ km/h los, aufeinander zu.
\begin{enumerate}[label=(\alph*),leftmargin=*,itemsep=3pt]
\item Nach welcher Zeit treffen sie sich? Gib die Zeit in Stunden und Minuten an.
\item Wie weit von Dorf A entfernt ist der Treffpunkt?
\item Nenne eine Annahme, die du für diese Rechnung machst und die in Wirklichkeit nicht ganz stimmt.
\end{enumerate}
\lpplatz{7}
\end{lpaufgabe}

\begin{lpaufgabe}{G8}{Einholen}{3 P}
Ein Lastwagen fährt mit $60$ km/h los. $30$ min später fährt ein Auto auf derselben Strasse hinterher, mit $90$ km/h.
\begin{enumerate}[label=(\alph*),leftmargin=*,itemsep=3pt]
\item Wie gross ist der Vorsprung des Lastwagens, wenn das Auto startet?
\item Nach welcher Zeit (ab dem Start des Autos) holt das Auto den Lastwagen ein?
\item Wie weit vom Start entfernt ist das?
\end{enumerate}
\lpplatz{5}
\end{lpaufgabe}
\end{document}
